\section{Introducción}

El presente plan de proyecto tiene como objetivo establecer las bases organizativas y metodológicas para el desarrollo del addon \textit{Análisis 3D CSV} en Blender. Este documento define la planificación del desarrollo, los recursos necesarios, los criterios de viabilidad y las estrategias de seguimiento del proyecto, con el fin de garantizar que se cumplan los objetivos establecidos en el TFG.

La definición de requisitos se realizó previamente al inicio del desarrollo. A partir de ese punto, el proyecto se ha desarrollado siguiendo un enfoque iterativo basado en ciclos de trabajo de dos semanas (sprints), permitiendo evolucionar el producto de forma incremental mediante distintas versiones funcionales del addon.

El desarrollo se ha basado en principios de metodologías ágiles, inspiradas en Scrum, mediante iteraciones cortas y evolución incremental del producto.

El repositorio GitHub ha actuado como herramienta central de gestión del proyecto, permitiendo el control de versiones, la trazabilidad del desarrollo y el seguimiento de la evolución del sistema mediante commits, incidencias (\textit{issues}) y versiones etiquetadas.

\section{Planificación temporal}

La planificación del proyecto se ha estructurado en sprints quincenales, alineados con la evolución real del desarrollo. El seguimiento de cada sprint se ha realizado mediante el repositorio de GitHub, utilizando \textit{commits}, \textit{issues} y versiones etiquetadas (\textit{tags}).

\begin{itemize}
    \item \textbf{Fase previa (Noviembre)}: Estudio de viabilidad del proyecto. Se realizaron pruebas iniciales para evaluar la posibilidad de desarrollar un sistema de recopilación y análisis de datos en Blender, incluyendo el análisis de su API, la viabilidad técnica de la captura de eventos y la posibilidad de generar métricas a partir de los datos registrados.
    \item 
    \item \textbf{Sprint 0 (Enero)}: Definición de requisitos, análisis del problema y estudio de la API de Blender.

    \item \textbf{Sprint 1 (Febrero - semanas 1-2)}: Desarrollo del primer addon funcional. Implementación inicial de la captura de eventos y exportación a CSV.
    \textit{Resultado: versión inicial funcional del sistema (tag: v0.1)}

    \item \textbf{Sprint 2 (Febrero - semanas 3-4)}: Mejora del primer addon, refinando la captura de datos y estructuración de la información.  

    \item \textbf{Sprint 3 (Marzo - semanas 1-2)}: Desarrollo del segundo addon, ampliando funcionalidades y mejorando la lógica de procesamiento de datos.
    \textit{Resultado: ampliación del sistema (tag: v0.2)}

    \item \textbf{Sprint 4 (Marzo - semanas 3-4)}: Incorporación de nuevas métricas y mejora del análisis de datos.  
    

    \item \textbf{Sprint 5 (Abril - semanas 1-2)}: Refactorización del sistema y mejora significativa de funcionalidades clave.  
    \textit{Commit destacado: "Refactorización del sistema de logging y detección de eventos (v0.4)"}  

    En este sprint se introducen mejoras relevantes:
    \begin{itemize}
        \item Soporte completo para Edit Mode mediante \textit{bmesh}.
        \item Mejora en la detección de operadores (Shift+D, Alt+D, Ctrl+V, Merge).
        \item Detección robusta de coordenadas UV.
        \item Optimización del sistema de logging, registrando únicamente cambios reales.
        \item Implementación de controles de ejecución (Start/Stop).
        \item Incorporación de un panel de depuración.
        \item Mejora en la generación y estructura de los archivos CSV.
    \end{itemize}

    \textit{Resultado: versión avanzada y estable del sistema (tag: v0.4)}

    \item \textbf{Sprint 6 (Abril - semanas 3-4)}: Mejora de la usabilidad y ajustes funcionales adicionales.  
    

    \item \textbf{Sprint 7 (Mayo - semanas 1-2)}: Pruebas funcionales y corrección de errores.  
     

    \item \textbf{Sprint 8 (Mayo - semanas 3-4)}: Documentación final y preparación de la entrega.  
    \textit{Resultado: versión final del sistema (tag: v1.0)}
\end{itemize}

La información reflejada en esta planificación se corresponde con la evolución real del repositorio del proyecto en GitHub, donde pueden consultarse los commits y versiones asociadas a cada iteración.

Este enfoque permite reflejar un desarrollo incremental del sistema, donde cada sprint aporta mejoras o nuevas funcionalidades, apoyándose en el control de versiones para garantizar la trazabilidad del proyecto.

\subsection{Resumen de sprints}

A continuación, se presenta un resumen de los sprints realizados durante el desarrollo del proyecto, incluyendo las principales tareas llevadas a cabo y los resultados obtenidos en cada iteración.

\begin{center}
\begin{tabular}{|c|c|p{5cm}|p{4cm}|}
\hline
\textbf{Sprint} & \textbf{Periodo} & \textbf{Funcionalidades principales} & \textbf{Resultado} \\
\hline

Sprint 1 & Febrero (1-2) & Estructura inicial del addon, captura de eventos y exportación a CSV & Versión inicial (v0.1) \\
\hline

Sprint 2 & Febrero (3-4) & Mejora del sistema de captura de datos y formato CSV & Sistema más robusto \\
\hline

Sprint 3 & Marzo (1-2) & Desarrollo del segundo addon y ampliación del procesamiento de datos & Versión ampliada (v0.2) \\
\hline

Sprint 4 & Marzo (3-4) & Implementación de métricas y análisis de datos & Sistema analítico inicial \\
\hline

Sprint 5 & Abril (1-2) & Refactorización del sistema, mejoras en logging y detección de eventos (Edit Mode, operadores, UV) & Versión avanzada (v0.4) \\
\hline

Sprint 6 & Abril (3-4) & Mejora de usabilidad, integración y optimización general & Sistema optimizado \\
\hline

Sprint 7 & Mayo (1-2) & Pruebas funcionales, validación y corrección de errores & Sistema validado \\
\hline

Sprint 8 & Mayo (3-4) & Documentación final y preparación de la entrega & Versión final (v1.0) \\
\hline

\end{tabular}
\end{center}

La información reflejada en esta tabla se corresponde con la evolución del repositorio del proyecto en GitHub.

\subsection{Viabilidad económica}

Aunque el proyecto se desarrolla en un entorno académico, es posible estimar su coste aproximado considerando un escenario profesional. Para ello, se tienen en cuenta los costes de software, hardware y desarrollo.

\subsubsection{Costes de software}

El proyecto utiliza principalmente herramientas de software libre, por lo que no implica costes directos de licencia:

\begin{itemize}
    \item Blender (GPL): 0 €
    \item Python (PSF License): 0 €
    \item GitHub (plan gratuito): 0 €
\end{itemize}

No obstante, el uso de un sistema operativo propietario como Windows puede implicar un coste aproximado de licencia de 150 €.

\subsubsection{Costes de hardware}

Se ha utilizado un equipo personal con un valor estimado de 1200 €. Considerando una amortización en 4 años (48 meses) y una duración del proyecto de aproximadamente 4 meses, el coste imputable sería:

\[
\text{Coste hardware} = \frac{1200}{48} \times 4 = 100 \,€
\]

\subsubsection{Coste de desarrollo}

Para estimar el coste de desarrollo, se considera un perfil de desarrollador junior con un salario aproximado de 1200 € mensuales durante 4 meses:

\[
\text{Coste desarrollo} = 1200 \times 4 = 4800 \,€
\]

\subsubsection{Coste total estimado}

\[
\text{Coste total} = 150 + 100 + 4800 = 5050 \,€
\]

Por tanto, aunque el coste real del proyecto es reducido al tratarse de un entorno académico, su equivalente en un contexto profesional sería de aproximadamente 5050 €.

Este análisis demuestra que el sistema es económicamente viable, especialmente al basarse en herramientas de software libre que reducen significativamente los costes de desarrollo.

\section{Estudio de viabilidad}

\subsection{Viabilidad legal}

Dado que Blender se distribuye bajo licencia GPL, los addons desarrollados para esta plataforma deben ser compatibles con dicha licencia. Por ello, el producto final se distribuye bajo licencia GPL (GNU General Public License), garantizando la libertad de uso, modificación y distribución del software.

En caso de utilizar bibliotecas de terceros, estas deberán ser compatibles con dicha licencia (por ejemplo, MIT o BSD).

Asimismo, el sistema incluye mecanismos de consentimiento del usuario para la recopilación de datos, cumpliendo con la normativa vigente en materia de protección de datos.

\subsection{Análisis de riesgos}

Como principal riesgo del proyecto se identifica la complejidad de la API de Blender, especialmente en la detección de eventos en distintos modos de trabajo como \textit{Object Mode} y \textit{Edit Mode}. 

Este riesgo se ha mitigado mediante un desarrollo iterativo, pruebas continuas y refactorizaciones progresivas del sistema, como la realizada en la versión v0.4.

Este planteamiento permite alinear la planificación teórica con la implementación real del sistema, garantizando coherencia entre la documentación y el desarrollo del proyecto.